\section{Discussion}
\label{sec:discussion}

\subsection{Why volume and effectiveness came apart}

The headline of Section~\ref{sec:results} is a dissociation: the LLM produced
$\PDvolumeRatioLlmEvo\times$ EvoMaster's test volume, covered every operation, generated
\RQIIItotalLlm{} edge-case scenarios against manual's \RQIIItotalManual{}, and still killed
half as many faults. Treating this as a paradox is a mistake. It follows directly from
\emph{what kind of assertion each arm writes}.

The LLM reads an OpenAPI specification. A specification documents \emph{status codes}: this
path returns 200, that one 400 on bad input. Tests synthesised from it therefore assert on
the status code and, at best, on the shape of the response body --- an oracle of the form
``the response is in the plausible class for this request.'' Our manual EP/BVA suite, drawn
from the same specification, writes the same kind of assertion. Both are
\emph{specification-derived status oracles}.

Now consider what our mutants do. A relational or arithmetic operator replacement inside a
numerical routine --- \texttt{<} becoming \texttt{<=}, \texttt{/} becoming \texttt{*} ---
changes a \emph{computed value}. The endpoint still accepts the request, still runs to
completion, still returns 200. Only the number in the body is now wrong. A status oracle
cannot see this. The test passes on the mutant exactly as it passed on the original, and
the fault survives.

EvoMaster writes a different kind of oracle. Its black-box mode records the actual response
of the original system and emits assertions pinning those recorded values. That is a
\emph{regression oracle}: it fails whenever any observed value changes, whether or not the
status does. Against value-changing mutants this is precisely the right instrument.

So the arms were never measuring the same thing. The LLM and manual arms answer ``does the
API respond in the documented class?''; EvoMaster answers ``does the API still compute what
it computed before?'' Our fault population consists almost entirely of value faults, so it
interrogates the second question. EvoMaster's win is not evidence that search-based
generation is smarter than an LLM --- it is evidence that a regression oracle beats a status
oracle on value faults, which is close to a tautology once stated.

Three observations support this account over alternatives.

\textbf{The per-subject pattern matches the prediction.} If oracle kind drives the result,
EvoMaster's advantage should concentrate where faults are most value-like. It does:
EvoMaster leads decisively on \texttt{rest-ncs} (\RQIINcsEvo{} against \RQIINcsLlm{}), whose
controllers are numerical routines where nearly every mutant shifts a computed value while
preserving the status. On \texttt{rest-scs}, which has more string and validation logic and
so more mutants that flip a response \emph{class}, the LLM leads (\RQIIScsLlm{} against
\RQIIScsEvo{}). The ordering inverts with fault character, exactly as the mechanism
predicts, and not with anything about the generators' sophistication.

\textbf{The LLM wins where its oracle is adequate.} On the documented error surface --- where
being wrong \emph{means} returning a different status --- the LLM leads every arm,
triggering \PDerrTriggeredLlm{} of \PDerrKey{} error behaviours against manual's
\PDerrTriggeredManual{} and EvoMaster's \PDerrTriggeredEvo{}, and eliciting a genuine 500
crash. Where a status oracle suffices, the LLM's superior volume converts into superior
detection. Where it does not, volume buys nothing.

\textbf{The degenerate LLM-manual tie is the mechanism's signature.} Two arms with zero
discordant pairs over \RQIIn{} faults is a startling result, and the same-author confound
(Section~\ref{sec:results-degenerate}) is a genuine part of the explanation. But it cannot
be the whole of it, because both suites are also \emph{structurally} blind to the same fault
class: any two status-oracle suites, however independently authored, are constrained to miss
every value fault. Independent authorship would move which \emph{status-visible} faults each
catches; it would not let either see a wrong number behind a 200. We therefore expect an
independently authored manual suite to remain close to the LLM on this fault population ---
a prediction our design leaves testable rather than resolves, and we flag it as such rather
than presenting the tie as a finding.

\subsection{Implications}

\textbf{For practitioners.} An LLM generator is a strong \emph{breadth} instrument and a
weak \emph{depth} one. It will reach every documented operation and probe the documented
error surface harder than a human or a search tool. It will not notice that your pricing
routine now returns the wrong total, because nothing in your specification says what the
right total is. The actionable rule: use LLM generation for specification conformance and
error-path breadth; do not let it displace regression oracles, property-based assertions, or
example-based expected values on computational paths. The two are complements, and our data
say the complement is not optional --- half the detected faults came from the arm the LLM
would have replaced.

\textbf{For the field.} Every metric in our study that the surveyed literature actually uses
--- coverage, scenario counts, 500-error tallies --- ranked the LLM first. The ground truth
ranked it second. That is not a quirk of our subjects; it is a structural consequence of
scoring test suites on volume proxies. When AutoRestTest reports 42 server errors against
EvoMaster's 20~\cite{autoresttest}, or LlamaRestTest 204 faults against 130~\cite{llamaresttest},
those numbers are real but they measure \emph{how hard the documented error surface was
hit}, not \emph{what fraction of faults was found}. Our results show the two can order
arms oppositely on the same code. Per-endpoint scenario counts --- the metric GAP-M asks
for and which we introduce --- are subject to the same caveat, and we would rather say so
than let our own contribution be misread: RQ3 is a measure of what the LLM \emph{writes},
and Section~\ref{sec:results-pvd} shows how little that predicted about what it
\emph{caught}.

\textbf{For benchmark design.} The field needs pre-seeded-fault benchmarks not because
Recall is a nicer statistic, but because without a denominator these two orderings are
indistinguishable. We would additionally argue that a benchmark should stratify its fault
population by \emph{oracle-observability} --- status-visible versus value-only --- since our
results show that single axis reordering the arms. A benchmark of only value faults makes
regression-oracle tools look dominant; one of only status faults makes LLMs look dominant.
Neither is a fact about the generators.

\subsection{Why the low absolute Recall is informative}

All three arms score low in absolute terms: the best, EvoMaster, kills \PDkilledEvo{} of
\PDkilledTotal{} mutants (\RQIIrecallEvoPct\%). We report this rather than tuning it away,
because it is a finding about black-box REST testing rather than a defect in our setup. Our
mutants live in core logic reachable only through the HTTP contract, and many are
\emph{equivalent-at-the-boundary}: the value they corrupt never surfaces in a response
field, or surfaces only for inputs no arm generated. A white-box tool with code-coverage
feedback would score higher. The honest reading is that black-box REST suites --- LLM,
human, or search-based --- detect a small minority of seeded logic faults, and the
interesting question is not which arm wins a low-scoring race but which fault classes all
three systematically cannot see.
